\problemname{Säker diktatorsortering}

Det är relativt välkänt att en korrekt sorteringsalgoritm inte kan använda färre än $O(N\log(N))$ jämförelser
i värsta fallet om algoritmen endast kan jämföra och byta plats på element.
För att komma runt detta har algoritmen \textit{diktatorsortering} utvecklats.
Den sorterar en lista på följande vis:

Behåll alltid det första elementet.
För varje efterföljande element $x$, jämför $x$ med det senast behållna elementet.
Om $x$ är \emph{strikt} större, behåll det.
Annars ta bort elementet $x$ från listan och skicka det till ett träningsläger.
Träningslägret är en lista där nya element läggs till längst bak.

Exempelvis: om listan är $[2,1,4,5,3,7]$, kommer den kvarvarande listan att vara $[2,4,5,7]$,
och träningslägret innehåller $[1,3]$.
Den här algoritmen går att implementera i $O(N)$ tid, men offrar korrekthet.

För att försöka fixa detta bestämmer du dig för att modifiera algoritmen.
Din idé är att köra diktatorsortering igen på listan i träningslägret.
Då får du ett nytt träningsläger att köra diktatorsortering på, och så vidare tills träningslägret blir tomt.
När algoritmen är klar har du flera listor som alla är sorterade; då är du färdig.
Det är högst oklart vad detta är bra för, men det känns som att du verkligen har bidragit till forskningen.
Därför bestämmer du dig för att döpa din nya algoritm till \textit{säker diktatorsortering}.

Exempelvis: anta att säker diktatorsortering körs på följande lista: $[5,3,1,4]$.
Efter första steget är resultatet listan $[5]$ och träningslägret $[3,1,4]$.
Om diktatorsortering sedan körs på träningslägret blir resultatet $[3,4]$ och det nya träningslägret $[1]$.
Eftersom $[1]$ är sorterad är algoritmen klar.
Det slutgiltiga resultatet blir då de tre listorna $[5]$, $[3,4]$ och $[1]$.

Tyvärr går säker diktatorsortering inte lika snabbt som vanlig diktatorsortering.
Kan du implementera algoritmen effektivt?
Det vill säga: efter varje gång diktatorsortering körs, skriv ut den resulterande sorterade listan,
tills algoritmen är färdig.


\section*{Indata}
Den första raden av indata innehåller heltalet $N$ ($1 \leq N \leq 2 \cdot 10^5$),
antalet tal i listan.

Den andra raden av indata innehåller $N$ heltal $a_1, \dots, a_N$ ($1 \leq a_i \leq 10^6$):
talen i listan som ska ``sorteras''.

\section*{Utdata}
Skriv ut innehållet i den sorterade listan på en egen rad efter varje gång vanlig diktatorsortering körs,
inklusive den sista.

\section*{Poängsättning}
Din lösning kommer att testas på en mängd testfallsgrupper.
För att få poäng för en grupp så måste du klara alla testfall i gruppen.

\noindent
\begin{tabular}{| l | l | p{12cm} |}
  \hline
  \textbf{Grupp} & \textbf{Poäng} & \textbf{Gränser} \\ \hline
  $1$    & $15$   & $N \leq 100$ \\ \hline
  $2$    & $30$   & $N \leq 5000$ \\ \hline
  $3$    & $10$   & $a_i \leq 2$ \\ \hline
  $4$    & $14$   & $a_i \leq 10$ \\ \hline
  $5$    & $31$   & Inga ytterligare begränsningar. \\ \hline
\end{tabular}

